\makeatletter
\ltx@ifpackageloaded{glossaries-extra}{

\renewcommand{\glossarypreamble}{\glsfindwidesttoplevelname[\currentglossary]}
\glsfindwidesttoplevelname
\setglossarystyle{alttree}
\setabbreviationstyle{short-long}
\makeglossaries

%---------------------- Batas Aman -----------------------%

%=========================================================%
%                      DAFTAR ISTILAH                     %
%                                                         %
%            Tambahkan istilah dengan memakai:            %
%                                                         %
%            \newglossaryentry{KATA_TUNJUK}{              %
%                name={NAMA_ISTILAH},                     %
%                description={PENJELASAN}                 %
%            }                                            %
%                                                         %
%   Istilah hanya akan muncul di halaman Daftar Istilah   %
%             jika disebut di dalam teksnya               %
%                                                         %
% Jangan lupa untuk menghapus contoh istilah di bawah ini %
%=========================================================%

\newglossaryentry{latex}{
    name={\LaTeX},
    description={Sistem penyiapan dokumen untuk tipografi berkualitas tinggi, yang sangat populer di kalangan akademisi untuk menulis dokumen ilmiah dan teknis}
}

\newglossaryentry{lualatex}{
    name={Lua\LaTeX},
    description={Pengembangan dari \LaTeX yang mengintegrasikan mesin pemrograman Lua, sehingga memberikan fleksibilitas tinggi dalam kustomisasi \textit{font} sistem secara dinamis}
}

\newglossaryentry{texmaker}{
    name={TeX Maker},
    description={Perangkat lunak penyunting (\textit{editor}) \LaTeX lintas \textit{platform} yang gratis dan terintegrasi, yang menyediakan antarmuka pengguna yang ramah untuk pemula}
}

\newglossaryentry{texstudio}{
    name={TeX Studio},
    description={Perangkat lunak penyunting (\textit{editor}) \LaTeX \textit{open-source} tingkat lanjut yang bercabang dari TeX Maker, yang dilengkapi fitur bantuan pengetesan kode secara interaktif dan pembacaan otomatis (\textit{code autocomplete})}
}

\newglossaryentry{miktex}{
    name={MiKTeX},
    description={Distribusi \LaTeX\ modern yang sangat populer untuk sistem operasi Windows, yang memiliki keunggulan berupa pengunduhan paket otomatis saat kompilasi berjalan}
}

\newglossaryentry{texlive}{
    name={TeX Live},
    description={Distribusi \LaTeX\ standar komprehensif lintas \textit{platform} (Linux, macOS, Windows) yang dirawat secara berkala oleh \textit{TeX Users Group}}
}

\newglossaryentry{microtype}{
    name=Microtype,
    description={Paket \LaTeX\ yang menangani optimasi tipografi tingkat mikro, seperti penyesuaian tonjolan huruf pada margin (\textit{protrusion}) dan kelebaran jarak huruf untuk meratakan teks secara sempurna}
}

%=========================================================%
%                    DAFTAR SINGKATAN                     %
%                                                         %
%           Tambahkan singkatan dengan memakai:           %
%                                                         %
%     \newacronym{KATA_TUNJUK}{SINGKATAN}{KEPANJANGAN}    %
%                                                         %
%  Singkatan hanya akan muncul di halaman Daftar Istilah  %
%             jika disebut di dalam teksnya               %
%                                                         %
%     Jangan lupa untuk menghapus contoh di bawah ini     %
%=========================================================%

\newacronym{wysiwyg}{WYSIWYG}{\textit{What You See Is What You Get}}
\newacronym{ctan}{CTAN}{\textit{Comprehensive TeX Archive Network}}
\newacronym{pdf}{PDF}{\textit{Portable Document Format}}
\newacronym{ide}{IDE}{\textit{Integrated Development Environment}}

%---------------------- Batas Aman -----------------------%

}{
    \typeout{Template Info: glossaries-extra not used.}
}
\makeatother